\section*{Problem 10}

We have 16 coins, one of which is odd (either heavier or lighter), and a two–pan
balance that only reports “balance’’ or “do not balance.’’ We do \emph{not} need
to know whether the odd coin is heavier or lighter, only which coin it is.

Each weighing can be viewed as a yes/no question of the form
“Is the odd coin among this set of coins?’’ so every weighing gives one bit of
information. With \(k\) weighings we can distinguish at most \(2^k\) possibilities.
Since there are 16 coins, we must have
\[
2^k \ge 16 \quad\Rightarrow\quad k \ge 4.
\]
So at least 4 weighings are necessary.

\medskip
\noindent\textbf{Strategy using 4 weighings.}

Label the coins \(1,2,\dots,16\) and write each label in 4-bit binary:

\[
\begin{array}{c|c}
\text{Coin} & \text{Binary label}\\\hline
1 & 0001\\
2 & 0010\\
3 & 0011\\
4 & 0100\\
\vdots & \vdots\\
16&1111
\end{array}
\]

For \(i=1,2,3,4\) (from least significant bit to most), perform weighing \(i\) as
the yes/no test:

\[
\text{Weighing }i:\quad
\text{``Is the odd coin among those whose $i$-th bit is $1$?''}
\]

Operationally, we put all coins with \(i\)-th bit \(1\) on the left pan and the
same number of other coins (chosen so that the total number of coins on each
side is the same) on the right pan. If the scale balances, the odd coin is not
in this subset; if it does not balance, then the odd coin is in this subset.

Record the 4-bit pattern of answers, using
\[
\text{``no''} \to 0,\qquad \text{``yes''} \to 1.
\]

For example, if the outcomes are
\[
\text{(yes, no, yes, yes)} = 1011_2,
\]
then the odd coin is the one with label \(1011_2 = 11\).

Every coin has a unique 4-bit label, and the 4 yes/no answers reproduce that
label, so after 4 weighings we can identify the odd coin uniquely.

\medskip
Thus 4 weighings are both necessary and sufficient.
